\documentclass[]{article}
\usepackage{xeCJK}        % XeLaTeX 中文支持
\usepackage{fontspec}     % 设置西文字体
\setCJKmainfont{SimSun}   % 设置中文主字体（宋体）
\setmainfont{Times New Roman} % 设置西文字体

% 数学公式包
\usepackage{amsmath, amssymb, amsfonts, amsthm}   % AMS数学包
\usepackage{mathtools}           % AMS扩展公式
\usepackage{bm}                   % 粗体数学符号

% ================= 排版美化 =================
\usepackage{geometry}             % 页边距
\geometry{left=2.5cm,right=2.5cm,top=2.5cm,bottom=2.5cm}
\usepackage{setspace}             % 行间距
\onehalfspacing                    % 1.5倍行距
\usepackage{titlesec}             % 调整章节标题
\usepackage{hyperref}             % 超链接
\hypersetup{
	colorlinks=true,
	linkcolor=blue,
	citecolor=cyan,
	urlcolor=magenta
}

%opening
\title{交换代数}
\author{today}
\author{Tamagawa}

\begin{document}
	
	

\maketitle

\section{Abstract}
这里笔者假定读者已经熟悉基本定义

\section{环与理想}

Proposition 1.2. Let $A$ be a ring $\neq 0$ . Then the following are equivalent: 

i) $A$ is a field; 

ii) the only ideals in $A$ are 0 and (1); 

iii) every homomorphism of $A$ into a non-zero ring $B$ is injective. 

p is prime $\Leftrightarrow A / \mathfrak{p}$ is an integral domain; 

m is maximal $\Leftrightarrow A / \mathfrak{m}$ is a field (by (1.1) and (1.2)). 

Theorem 1.3. Every ring $A \neq 0$ has at least one maximal ideal. (Remember that "ring" means commutative ring with 1.) 

Corollary 1.4. If $\mathfrak{a} \neq (1)$ is an ideal of $A$ , there exists a maximal ideal of $A$ containing $\mathfrak{a}$ . 

Corollary 1.5. Every non-unit of $A$ is contained in a maximal ideal. 

Proposition 1.6. i) Let $A$ be a ring and $\mathfrak{m} \neq (1)$ an ideal of $A$ such that every $x \in A - \mathfrak{m}$ is a unit in $A$ . Then $A$ is a local ring and $\mathfrak{m}$ its maximal ideal.

ii) Let $A$ be a ring and $\mathfrak{m}$ a maximal ideal of $A$ , such that every element of  $1 + m$ (i.e., every $1 + x$ , where $x \in m$ ) is a unit in $A$ . Then $A$ is a local ring. Proof. i) Every ideal $\neq (1)$ consists of non-units, hence is contained in $m$ . Hence $m$ is the only maximal ideal of $A$ . 

Proposition 1.7. The set $\mathfrak{N}$ of all nilpotent elements in a ring $A$ is an ideal, and $A / \mathfrak{N}$ has no nilpotent element $\neq 0$ . 

Proposition 1.8. The nilradical of $A$ is the intersection of all the prime ideals of $A$ . 

Proposition 1.9. $x \in \Re \Leftrightarrow 1 - xy$ is a unit in $A$ for all $y \in A$ . 

Again, in $\mathbf{Z}$ , we have $(a + b)(a \cap b) = ab$ ; but in general we have only $(a + b)(a \cap b) \subseteq ab$ (since $(a + b)(a \cap b) = a(a \cap b) + b(a \cap b) \subseteq ab$ ). Clearly $ab \subseteq a \cap b$ , hence $$
a \cap b = a b \text {   provided   } a + b = (1).
$$

Two ideals $a, b$ are said to be coprime (or comaximal) if $a + b = (1)$ . Thus for coprime ideals we have $a \cap b = ab$ . Clearly two ideals $a, b$ are coprime if and only if there exist $x \in a$ and $y \in b$ such that $x + y = 1$ . 

Proposition 1.10. i) If $a_{i}, a,$ are coprime whenever $i \neq j$ , then $\Pi a_{i} = \bigcap a_{i}$ . 

ii) $\phi$ is surjective $\Leftrightarrow \mathfrak{a}_i, \mathfrak{a}_j$ are coprime whenever $i \neq j$ . 

iii) $\phi$ is injective $\Leftrightarrow \bigcap a_{i} = (0)$ . 

Proposition 1.11. i) Let $\mathfrak{p}_1, \ldots, \mathfrak{p}_n$ be prime ideals and let $\mathfrak{a}$ be an ideal contained in $\bigcup_{i=1}^n \mathfrak{p}_i$ . Then $\mathfrak{a} \subseteq \mathfrak{p}_i$ for some $i$ . 

ii) Let $\mathfrak{a}_1, \ldots, \mathfrak{a}_n$ be ideals and let $\mathfrak{p}$ be a prime ideal containing $\bigcap_{i=1}^{n} \mathfrak{a}_i$ . 

Then $\mathfrak{p} \supseteq a_i$ for some $i$ . If $\mathfrak{p} = \bigcap a_i$ , then $\mathfrak{p} = a_i$ for some $i$ . 

If a, b are ideals in a ring A, their ideal quotient is $$
(a: b) = \{x \in A: x b \subseteq a \}
$$

Example. If $A = \mathbf{Z}$ , $a = (m)$ , $b = (n)$ , where say $m = \prod_{p} p^{\mu_p}$ , $n = \prod_{p} p^{\nu_p}$ , then $(a:b) = (q)$ where $q = \prod_{p} p^{\nu_p}$ and $$
\gamma_ {p} = \max (\mu_ {p} - \nu_ {p}, 0) = \mu_ {p} - \min (\mu_ {p}, \nu_ {p}).
$$

Hence $q = m / (m, n)$ , where $(m, n)$ is the h.c.f. of $m$ and $n$ . 

\end{document}
